% =============================================================================
% Part I -- The Constructive Reach
% =============================================================================
\noindent Part~I answers Clause~1: \textbf{what does the discrete ontology
force?} It exhibits the algebraic spine (the theorem-grade results the
substrate produces) in the order a mathematician would build them, each
object from prior ones, each step independently checkable. The construction runs
$i \to \Zi \to \Gs \to \text{master quadratic} \to \text{spine}$. It carries
\textbf{no physics identification}: the famous numerical match $x_+ = 1/\alpha$
is a \emph{match}, deferred to Part~III, and Part~I says so at the one place it
is tempting to overstep.

\medskip
\noindent The spine is \textbf{nine numbered results: seven are theorem-grade,
and two are honestly tiered below theorem grade} (the coefficient-16 identity,
whose structural necessity is conjectural, and the Phase-J ultralocality, which
is a theorem only at $L=2$). \Cref{sec:spine} closes by re-stating that count.

\section{The seed: \texorpdfstring{$i$}{i} and \texorpdfstring{$\Zi$}{Z[i]}}\label{sec:seed}
The construction begins with one structural choice and one group-theory fact,
and it is worth separating them cleanly, because the honesty of everything
downstream depends on not confusing the two.

\runin{The structural choice.}\margetag{ax} The substrate is the cubic lattice of
Part~0. Its coordinate planes carry an order-4 rotational symmetry: a $90^\circ$
rotation of the $(x,y)$ plane sends
$(1,0)\mapsto(0,1)\mapsto(-1,0)\mapsto(0,-1)\mapsto(1,0)$, a cycle of length~4.
FTD's seed move is to \emph{read this order-4 planar symmetry as the arithmetic
of the Gaussian integers} $\Zi=\{a+bi : a,b\in\Z\}$, with the $90^\circ$ rotation
realized as multiplication by $i$. That the lattice's quarter-turn \emph{should
be read through} $\Zi$ rather than through some other order-4 structure is a
modelling decision, not a theorem, the discrete-ontology analogue of choosing
to build $\C$ as $\R[i]$: natural, motivated, and a \emph{choice}. We mark it an
\etag{ax} and flag it as such, so that no later theorem silently inherits the
status of ``forced'' from a step that was actually ``chosen.''

\runin{The group-theory facts.}\margetag{thm} Once $\Zi$ is the substrate's
arithmetic, the following are ordinary, checkable theorems of algebra, owing
nothing further to FTD:
\begin{itemize}[leftmargin=1.4em, itemsep=3pt]
  \item The unit group of $\Zi$ is $\Ziu=\{1,i,-1,-i\}\cong\Z/4$, of order
        $|\Ziu|=4$. \emph{(Proof: $a+bi$ is a unit iff its norm $a^2+b^2=1$,
        whose only integer solutions are $(\pm1,0)$ and $(0,\pm1)$.)}
  \item $\Zi$ is a Euclidean domain, hence a PID and a UFD; its primes split
        into the three Gauss classes: the ramified prime $1+i$ (norm~2), the
        split rational primes $p\equiv 1\ (\mathrm{mod}\ 4)$, and the inert
        rational primes $p\equiv 3\ (\mathrm{mod}\ 4)$ (Fermat's two-square
        theorem).
\end{itemize}
The single integer this seed deposits into the rest of the construction is
$4=|\Ziu|$. It reappears, squared, as the coefficient~16 of the master quadratic
(\secref{sec:masterquad}), and the split/inert residue classes mod~4 reappear as
the arithmetic content of one of the four $\Gs$ constructions (\secref{sec:gstar}).
Hold the distinction firmly: $4$ is a \emph{theorem} about $\Zi$; that the
\emph{lattice} is entitled to $\Zi$ in the first place is the axiom.

\section{\texorpdfstring{$\Gs$}{G*} in four representations of differing provenance}\label{sec:gstar}
The bridge constant is\margetag{thm}
\[ \Gstardef. \]
What earns $\Gs$ its place at the centre of the spine is that it admits four
\emph{representations} of sharply differing character: a $\Gamma$-ratio, a
lattice Green's-function period, a regularized determinant ratio, and a finite
product. These are not four independent witnesses: three share the quarter-integer
data, and the Watson value is itself a $\Gamma(1/4)$-period, so all four are
representations of one and the same CM period of the lemniscatic curve. What is
striking is that a single period wears four such different faces (arithmetic,
lattice-analytic, spectral, and finitary), one of which (Watson's
body-centred-cubic integral) is of genuinely independent provenance. Each is shown
where it is cheap and cited where it is deep.

\runin{Route 1: the $\Gamma$-ratio.} Directly from the definition and the
Euler reflection formula,
\[ \reflection, \]
multiplying $\Gs=\Gamma(1/4)/\Gamma(3/4)$ by $\Gamma(3/4)^2$ and substituting
gives the two standard closed forms
\[ \Gstarclosed, \]
where $\varpidef$ is the Bernoulli/Gauss lemniscate constant. A three-line proof,
verified to 50 digits.

\runin{Route 2: the Watson BCC period}, conditional on Watson 1939. Let $W_3$
be the body-centred-cubic Watson integral
\[ \watsonint. \]
Watson's 1939 closed-form evaluation~\cite{watson1939}, consolidated by
Glasser--Zucker~\cite{glasserZucker1980}, gives
\[ \watsonval. \]
The BCC sub-lattice is one of the three polyhedral layers of the substrate's
Moore neighbourhood, so this is $\Gs$ arising from a \emph{lattice}
quantity, a Green's-function period of an entirely different character from
Route~1's $\Gamma$-ratio. ``Conditional on Watson 1939'' means the rigour equals
that of the cited classical evaluation, not that it rests on a conjecture.
Verified to 100 digits in PARI.

\runin{Route 3: the $\det_\zeta$ quarter-conjugacy bridge.} Let $J$ be the
conjugacy operator with $J^2=-I$ (multiplication by $i$; the same $Z_4$ seed of
\secref{sec:seed}). A wavefunction on $S^1$ with the quarter-twisted boundary
condition $\psi(\varphi+2\pi)=J\,\psi(\varphi)$, evolved by the natural
quarter-period operator $-i\,d/d\varphi$, carries its two $J$-eigensectors on the
shifted spectra $D_{1/4}=\{n+1/4\}_{n\ge0}$ and $D_{3/4}=\{n+3/4\}_{n\ge0}$ (the
choice of that operator is a modelling choice, like the seed, not something
forced). Lerch's formula for the $\zeta$-regularized
determinant of an arithmetic progression,
\[ \lerch, \]
yields the determinant ratio
\[ \detzratio, \]
the $\sqrt{2\pi}$ cancelling. The arithmetic is the reason this route belongs to
the same story as the seed: $4\cdot D_{1/4}=\{n\equiv1\ (\mathrm{mod}\ 4)\}$ and
$4\cdot D_{3/4}=\{n\equiv3\ (\mathrm{mod}\ 4)\}$ are exactly the two non-trivial
residue classes mod~4, and, restricted to primes, the split and inert prime
classes of $\Zi$ (\secref{sec:seed}). $\Gs$ is the regularized asymmetry between
them.

\runin{Route 4: the finite-$N$ attractor.} Define the finite product
\[ \finiteN. \]
Every $G^{*}_N$ is computable in $O(N)$ rational operations with no transcendental
ever invoked ($G^{*}_0=3$, $G^{*}_1=21/(5\sqrt2)\approx2.9698$,
$G^{*}_5\approx2.95995$, $G^{*}_{20}\approx2.95878$). A Stirling expansion of
$\Gamma(N+7/4)/\Gamma(N+5/4)\sim(N+1)^{1/2}$ gives
\[ \finiteNrate \]
(empirical $C\approx0.046$). This route does double duty. It is a fourth
construction of $\Gs$, and it \textbf{discharges the Part-0
$\varepsilon$--$L$ obligation} for $\Gs$: the undefined-boundary ontology forbids
treating $\Gs=\lim_{N\to\infty}G^{*}_N$ as a primitive completed limit, and
Route~4 supplies exactly the finitary restatement the ontology demands: $G^{*}_N$
is defined for every finite $N$ and approaches $\Gs$ at the stated rate. $\Gs$ is
therefore reachable by a convergent finite computation, no completed infinity
required.

\runin{$\Gs$ is the lemniscatic ratio, not $\varpi$.} $\Gs\approx\Gstarapprox$ is
the ratio $\Gamma(1/4)/\Gamma(3/4)$. It is \textbf{not} the Bernoulli/Gauss
lemniscate constant $\varpi\approx\varpiapprox$, a distinct number related by
$\Gs=2\varpi/\sqrt\pi$. Informal usage sometimes calls both ``the lemniscate
constant''; here they are never conflated, and the distinction is load-bearing in
the most literal sense. The master quadratic of \secref{sec:masterquad} produces
$x_+=\xplusnum$ \textbf{only} at $\Gs=\Gstarapprox$; substituting
$\varpi=\varpiapprox$ into the same polynomial gives $x_+\approx107.3$, nowhere
near $1/\alpha$.

\section{The master quadratic}\label{sec:masterquad}
With $\Gs$ in hand, the construction's central polynomial is purely algebraic.

\begin{forcedbox}[The master quadratic \etag{thm}]
Define
\[ \masterquad. \]
Its discriminant factors cleanly,
\[ \discriminant \]
which is \textbf{positive} because $\Gs>1/4$ (indeed $\Gs\approx2.96$), so both
roots are real:
\[ \rootsexplicit \]
Numerically $x_+=\xplusnum$ and $x_-=\xminusnum$.
\end{forcedbox}

\noindent This is the quadratic formula applied to a polynomial whose
coefficients are integer multiples of powers of $\Gs$; nothing more is involved.
Verified numerically across 54 independent checks.

\runin{The coefficient-16 soft spot, flagged honestly.}\margetag{cnj} The
prefactor~16 has a tempting arithmetic reading: for the lemniscatic CM elliptic
curve $E:y^2=x^3-x$, the automorphism group over $\overline{\Q}$ is
$\Aut(E)=\{\pm1,\pm i\}\cong\Z/4$ (the unit group of $\Zi$ again), so
$|\Aut(E)|=4$ and $|\Aut(E)|^2=16$. That $|\Aut(E)|^2=16$ is a \etag{thm}. But
that the master quadratic's prefactor is \emph{forced} to equal $|\Aut(E)|^2$ (that
two objects equal 16 \emph{for a structural reason}) is not proven: it is a
\etag{cnj} on the structural necessity, the framework's softest mathematical
spot. A partial unification exists (a tower-level reading supplies a reason for
$k=4$, hence $2^k=16$, via $|\Ziu|^2=16$), but a partial unification is not a
forcing theorem, and the honest grade stays \etag{cnj} on necessity.

\runin{No physics here.} $x_+=137.036\ldots$ is, \emph{in Part~I}, a root of a
polynomial, nothing more. The proximity to $1/\alpha$ is the motivation for
the whole enterprise, but the identification $x_+ \leftrightarrow 1/\alpha$ is a
\etag{smc} deferred in full to \textbf{Part~III}; it is not asserted as forced
anywhere in the constructive Parts. The smaller root $x_-\approx\xminusapprox$ is
a pure algebraic artifact, pinned by the Vieta relation $x_-=16\Gsss/x_+$ the
instant $x_+$ is, and it carries no physics. The historical identification
$x_-\leftrightarrow N_c$ is \textbf{retired}; FTD's $N_c=3$ is independently
sourced from the Moore-neighbourhood topology, owing nothing to the polynomial's
smaller root.

\section{The spine built on top}\label{sec:spine}
On the foundation of $\Gs$ (\secref{sec:gstar}) and the master quadratic
(\secref{sec:masterquad}), a further family of theorem-grade results stands. Each
is stated at its exact grade and condition; short proofs are shown, deep external
dependencies are cited.

\runin{Harmonic-invariant tower.}\margetag{thm} For each integer $k\ge3$, the
$(1+i)$-tower master quadratic is $\towerquad$ (the $k=4$ instance is
\secref{sec:masterquad}'s polynomial). Writing the normalized roots
$y_\pm := x_\pm/\Gs$, at \emph{every} level $k\ge3$,
\[ \harmonicinv. \]
The proof is a three-line Vieta computation:
$1/x_+ + 1/x_- = (x_+ + x_-)/(x_+ x_-)
= (2^k\Gpow{k-2})/(2^k\Gpow{k-1}) = 1/\Gs$; multiplying through by $\Gs$ gives
$1/y_+ + 1/y_- = 1$. Verified to 50 digits.

\runin{CM-curve uniqueness}, under the trivial-multiplier criterion.\margetag{thm}
Among all imaginary quadratic fields $K=\Q(\sqrt{-d})$, the field $\Q(i)$ ($d=1$,
fundamental discriminant $-4$) is the \textbf{unique} one satisfying
$|\mu_K|=|\mathrm{disc}(K)|$: for $\Q(i)$, $|\mu_K|=4$ and
$|\mathrm{disc}(K)|=4$; for every other imaginary quadratic field the two differ.
\emph{(Proof sketch: $|\mu_K|\in\{2,4,6\}$ with 4 only at $d=1$ and 6 only at
$d=3$; $|\mathrm{disc}(K)|\ge3$ always, $=3$ only at $d=3$, $=4$ only at $d=1$;
case-checking the three unit-group orders leaves $d=1$ as the sole coincidence.)}
The trivial-multiplier criterion is load-bearing: a ``match'' here requires the
natural root to equal the target directly. Under a looser rational-multiplier
criterion, 20 additional non-canonical matches exist in the tested grid, and that
reading is a \etag{sel}, not a theorem. The proof is a complete case analysis over
all imaginary quadratic fields; the $d\le200$ computation is a redundant
cross-check.

\runin{$\Q(\Gs)$ is $\pi$-free}, conditional on Chudnovsky
1976.\margetag{thm} $\Q(\Gs)\cap\Q(\pi)=\Q$: the field generated by $\Gs$ shares
no algebraic content with $\Q(\pi)$ beyond the rationals. The proof reduces a
hypothetical polynomial relation between $\Gs$ and $\pi$ (via the identity
$\Gs\cdot\pi\cdot\sqrt2=\Gamma(1/4)^2$) to a polynomial relation between
$\pi$ and $\Gamma(1/4)$, which Chudnovsky's 1976 algebraic-independence
theorem~\cite{chudnovsky1984} forbids. ``Conditional'' here means the rigour
equals that of this published, foundational transcendence result (consolidated in
Waldschmidt~\cite{waldschmidt2000}), not that it rests on a conjecture.

\runin{Tower-discriminant transcendence}, conditional on
Schneider--Chudnovsky~\cite{schneider1937,chudnovsky1984}.\margetag{thm} For the
tower of the harmonic invariant, the discriminant factors as $\towerdisc$. Then
$A_k$ is rational at $k=3$ ($A_3=1$) and \textbf{transcendental over $\Q$} (i.e.\
$A_k\notin\overline{\Q}$) at every $k\ge4$ ($A_4=4\Gs-1$, etc.): a non-rational
polynomial with rational coefficients in the transcendental $\Gs$ takes
transcendental values.

\runin{BCC complex structure.}\margetag{thm} The 8 BCC corners under $90^\circ$
rotation form two orbits of size~4, and, as a module over the cyclic group
$C_4=\langle 90^\circ\text{ rotation}\rangle\cong\Z/4$,
\[ \bccdecomp, \]
where $V_{\mathrm{complex}}$ carries a natural $\Zi$-module structure
$\cong\Zi^2$. Paired with it is a clean no-go: there is \textbf{no} injective
homomorphism $\Ziu\to O_h^{\mathrm{ab}}$, since $\Ziu\cong\Z/4$ has an order-4
element but $O_h^{\mathrm{ab}}\cong\Z/2\times\Z/2$ (Klein four) does
not, a one-line group-order argument. Together these say the substrate's
$\Zi$ structure lives genuinely in the BCC layer but does \textbf{not} extend to a
global lattice-symmetry embedding. Verified in exact rationals.

\runin{Lemniscatic $L$-value}, conditional on Rubin 1991.\margetag{thm} The
central $L$-value of the lemniscatic curve $E_{\mathrm{lemn}}:y^2=x^3-x$ (Cremona
32.a3) has the clean closed form
\[ \Lvalueeq, \]
via the full BSD formula for the CM rank-0 case~\cite{rubin1991}. (An earlier
$\varpi/2$ came from a BSD-convention double-count; the corrected factor is
$\varpi/4$, verified to 27 digits against the curve's entry in
LMFDB~\cite{lmfdb}.)

\runin{Three further results, conditional on the classical theorems named.}%
\margetag{thm}
The character $\chi_{-4}$ four-level unification\,---\,$\chi_{-4}$ on
$(\Z/4\Z)^\times$ generating the entire $\Gs/\GG$ identity algebra through four
functorial projections (lattice, Chowla--Selberg, Hecke, Dirichlet)\,---\,holds
conditional on Deligne's period conjecture~\cite{deligne1979} in the CM case,
itself proved unconditionally (Blasius/Anderson/Shimura). The $\eta$-tower across
the nine class-number-one imaginary-quadratic fields,
\[ |\eta(\tau_K)|^{2w_K} \;=\; \frac{G_K^{\,w_K}}{(2\pi|d_K|)^{w_K/2}}, \]
holds conditional on Chowla--Selberg~\cite{chowlaSelberg1967}, unifying the
Heegner near-integer phenomenon as a $\chi$-projection. Here $G_K$ is the
field-specific Chowla--Selberg period (it equals $\Gs$ only for the Gaussian
field $d=1$), so the nine fields share the \emph{form} of the identity, not
the constant. The
$\mathrm{Sym}^2\oplus\mathrm{Sym}^3$ structure singles out the construction's
pair. Among the period polynomials $x^2-16\Gpow{a}x+16\Gpow{b}$ with $a<b$
positive integers and positive discriminant (which requires only
$b<2a+\log 4/\log\Gs\approx 2a+1.28$, i.e.\ $b\le 2a+1$ for integers), the pair
$(a,b)=(2,3)$ is distinguished by the \emph{parity of $\Gs$ beneath its
discriminant's radical: $\Delta_{(2,3)}=64\Gsss(4\Gs-1)$ carries an odd power of
$\Gs$}, so $\sqrt{\Delta}=8\Gs\,\sqrtwall$ introduces the generator
$\sqrtwall$ (the very square root that walls Part~II), whereas the
neighbouring pairs carry an even power ($(1,3)$ gives
$\sqrt{\Delta}\propto\sqrt{4-\Gs}$; $(1,2)$ and $(2,4)$ have $b=2a$). Equivalently,
the squarefree radicand of $\Delta$ retains a bare factor of $\Gs$ exactly on the
line $b=2a-1$ (the pairs $(2,3),(3,5),(4,7),\dots$, each adjoining the wall
field $\Q(\Gs)(\sqrtwall)$), of which $(2,3)$ is the minimal member with $a<b$.
That much is elementary algebra; the stronger reading (that $(2,3)$ is the
\emph{unique} natural target) is a \etag{sel}\margetag{sel}, argued from
naturalness, not asserted here as a theorem.

\runin{Phase-G geometric Coulomb.}\margetag{thm} The remaining theorem-grade
node is a lattice fact. On the cubic lattice the static potential between two
manifested sites solves the lattice Poisson equation, and its solution is, at
every finite $L$, exactly the periodic-lattice Poisson Green's function
$G_L$. The emergent inverse-square Coulomb tail is therefore lattice geometry,
carrying no fine-structure content of its own: the lattice coupling reads as a
zero-parameter geometric object, $\alpha_r(r,L)=2\,r\,G_L(r)$, not a free
constant. This is a finite-$L$ identity, established directly, with no continuum
limit required.

\runin{The two honestly-tiered results.} Two of the nine numbered results sit
\emph{below} theorem grade and are stated as such:
\begin{itemize}[leftmargin=1.4em, itemsep=4pt]
  \item \textbf{Coefficient-16}: the value-level identity $16=|\Aut(E)|^2$
        is true, but its structural necessity is a \etag{cnj}, as set out in
        \secref{sec:masterquad}.\margetag{cnj}
  \item \textbf{Phase-J partition-function ultralocality}: a
        \etag[ at $L=2$ only]{thm}.\margetag[ at $L=2$ only]{thm} On a $2^3$
        lattice the classical Euclidean action depends on the state field solely
        through $\sum_i s_i^2$ (the manifested-site count) and is invariant under
        spatial permutation of charge placement; this is proved by explicit
        construction. The general-$L$ extension is \textbf{disconfirmed}: the
        $L=2$ ultralocality is a Nyquist-mode counting degeneracy (every non-zero
        momentum has $\sin(k_i)=0$), not a structural property, and at $L\ge3$ the
        action does depend on spatial placement. The $L=2$ restriction is stated
        plainly and is the spine's honest claim.
\end{itemize}

\runin{The canonical accounting.} That completes the spine: \textbf{nine numbered
results, seven theorem-grade} ($\Gs$ identity, master quadratic, CM uniqueness,
Watson identity, Phase-G geometric Coulomb, harmonic-invariant tower, $\Q(\Gs)$
$\pi$-freeness conditional on Chudnovsky) \textbf{plus two tiered below theorem grade}
(coefficient-16 and Phase-J $L=2$). This is the discrete ontology's firm
ground (its answer to Clause~1), and it stands independent of any physics
interpretation. The match to physics is a separate matter, quarantined in
Part~III.

\section{The construction map}\label{sec:map}
% =============================================================================
% dag.tex  --  The construction DAG (Figure in Part I, section "The construction map")
% -----------------------------------------------------------------------------
% Hand-authored, fully vector, monochrome.  Differentiation is by SHAPE + RULE
% (never colour), mirroring the epistemic-tag tiers of the body:
%   double circle = AXIOM seed   diamond (heavy) = the G* hub
%   solid rounded box = theorem-grade   dashed rounded box = below-theorem
%   small square = framework integer    tiny box = (omitted; ledger anchors live in text)
%   solid arrow = "constructed from"     dashed arrow = "necessity conjectured / weak link"
%
% This figure is hand-set so it stays monochrome, legible, and faithful to the
% nine-result spine accounting (seven theorem-grade, two tiered).
% =============================================================================
\begin{figure}[t!]
  \centering
  \resizebox{\textwidth}{!}{%
  \begin{tikzpicture}[
      >={Stealth[length=4pt]},
      seed/.style   ={circle, double, draw, line width=0.7pt, inner sep=2.2pt,
                      font=\footnotesize, align=center},
      intg/.style   ={rectangle, draw, line width=0.5pt, inner sep=3pt,
                      font=\footnotesize, align=center},
      hub/.style    ={diamond, draw, line width=1.2pt, inner sep=2pt, aspect=1.6,
                      minimum width=1.25cm, minimum height=0.95cm,
                      font=\footnotesize\bfseries, align=center},
      thm/.style    ={rectangle, rounded corners=2pt, draw, line width=0.7pt,
                      inner sep=3pt, font=\footnotesize, align=center, text width=2.3cm},
      tier/.style   ={rectangle, rounded corners=2pt, draw, line width=0.5pt,
                      dash pattern=on 2.4pt off 1.8pt, inner sep=3pt,
                      font=\footnotesize, align=center, text width=2.3cm},
      con/.style    ={->, draw, line width=0.6pt},
      wk/.style     ={->, draw, line width=0.5pt, dash pattern=on 2.4pt off 1.8pt,
                      gray!60!black},
      lbl/.style    ={font=\scriptsize\itshape, fill=white, inner sep=1.2pt},
    ]
    % ---- main trunk (axiom -> integer -> hub -> quadratic -> tower) --------
    \node[seed] (zi)  at (4.0, 0.0)  {$\mathbb{Z}[i]$\\[-1pt]\tiny\scshape axiom};
    \node[intg] (four) at (4.0,-1.9) {$4=|\mathbb{Z}[i]^{\times}|$};
    \node[hub]  (gs)  at (4.0,-3.9)  {$G^{*}$};
    \node[thm]  (mq)  at (4.0,-6.9)  {Master quadratic\\ $x^2-16G^{*2}x+16G^{*3}$\\ \textsc{(t2)}};
    \node[thm]  (tow) at (4.0,-9.4)  {Harmonic\\ invariant tower \textsc{(t8)}};

    % ---- right branches off the hub G* ------------------------------------
    \node[thm]  (wat) at (8.2,-3.1)  {Watson identity\\ $W_3=G^{*2}/2\pi$ \textsc{(t5)}};
    \node[thm]  (qg)  at (8.2,-4.9)  {$\mathbb{Q}(G^{*})$ is\\ $\pi$-free \textsc{(t9)}};
    \node[thm]  (cm)  at (8.2,-6.7)  {CM-curve\\ uniqueness \textsc{(t3)}};
    \node[thm]  (bcc) at (8.2,-8.5)  {BCC complex\\ structure};

    % ---- left / lower: the two below-theorem-grade results (dashed) -------
    \node[tier] (c16) at (0.1,-6.9)  {Coefficient\\ $16=4^2$ \textsc{(t4)}\\ \tiny necessity conj.};
    \node[tier] (pj)  at (0.1,-9.4)  {Phase-J ultra-\\ locality \textsc{(t7)}\\ \tiny $L=2$ only};
    \node[thm]  (pg)  at (8.2,-10.3) {Phase-G geometric\\ Coulomb \textsc{(t6)}};

    % ---- edges -------------------------------------------------------------
    \draw[con] (zi)  -- (four);
    \draw[con] (four) -- (gs);
    \draw[con] (gs)  -- (mq);
    \draw[con] (mq)  -- (tow);
    % the "16 = 4^2" feed from the integer into the quadratic (curved, left)
    \draw[con] (four) to[out=200,in=150] node[lbl,pos=0.62]{$16=4^2$} (mq.west);
    % branches off G* -- the two lowest (bcc, pg) route via an eastward waypoint
    % so they descend on the RIGHT of the central column, clear of the mq box.
    \draw[con] (gs) to[out=18,in=180]  (wat.west);
    \draw[con] (gs) to[out=-8,in=170]  (qg.west);
    \draw[con] (gs) to[out=-30,in=170] (cm.west);
    \draw[con] (gs) to[out=-42,in=120] (6.5,-7.5) to[out=-60,in=170] (bcc.west);
    \draw[con] (gs) to[out=-52,in=120] (6.3,-9.0) to[out=-60,in=160] (pg.west);
    % weak link: 16 = 4^2 necessity is conjectured
    \draw[wk] (mq.west) to[out=180,in=0] (c16.east);
    % Phase-J hangs off the lattice axiom (independent ultralocality fact)
    \draw[wk] (zi.west) to[out=180,in=90] (pj.north);
  \end{tikzpicture}}
  \caption[The construction DAG]{\textbf{The construction DAG.}\enspace Shape
    and rule encode epistemic status: the double-ringed circle is the lone
    \textsc{axiom} ($\mathbb{Z}[i]$), the heavy diamond is the hub ($G^{*}$),
    solid boxes are theorem-grade, and dashed boxes are tiered below theorem
    grade (the coefficient-16 necessity, conjectured; Phase-J, a theorem only at
    $L=2$). The chain runs $i \to \mathbb{Z}[i] \to 4 \to G^{*} \to{}$ master
    quadratic $(16=4^{2}) \to{}$ spine; reading the arrows backward traces every
    theorem-grade node to the one structural choice at the root.}
  \label{fig:dag}
\end{figure}

The spine is not a list but a dependency structure: a directed acyclic graph
rooted in the seed. The axiom reading of the lattice's quarter-turn as $\Zi$
(\secref{sec:seed}) deposits the integer~4; $\Gs$ is constructed four ways from
quarter-integer data (\secref{sec:gstar}); the master quadratic $\masterquad$ is
assembled from $\Gs$ and $16=4^2$ (\secref{sec:masterquad}); and the spine
theorems (harmonic tower, CM uniqueness, $\Q(\Gs)$ $\pi$-freeness, the BCC
complex structure, the $L$-value and its companions) hang off $\Gs$ and the
quadratic in turn (\secref{sec:spine}). Reading the arrows backward is the audit
trail: every theorem-grade node traces, through explicit and checkable steps, to
the one structural choice at the root, and to nothing else.
